\documentclass[11pt]{article}

\usepackage[T1]{fontenc}
\usepackage{lmodern}
\usepackage[margin=1in]{geometry}
\usepackage{amsmath,amssymb,amsthm,mathtools}
\usepackage{enumitem}
\usepackage{booktabs,tabularx,array}
\usepackage{hyperref}
\usepackage{microtype}
\usepackage{xcolor}

\hypersetup{
  colorlinks=true,
  linkcolor=blue!50!black,
  citecolor=blue!50!black,
  urlcolor=blue!50!black
}

\newtheorem{theorem}{Theorem}[section]
\newtheorem{proposition}[theorem]{Proposition}
\newtheorem{corollary}[theorem]{Corollary}
\newtheorem{lemma}[theorem]{Lemma}
\newtheorem{definition}[theorem]{Definition}
\newtheorem{remark}[theorem]{Remark}
\newtheorem{placeholder}[theorem]{Placeholder}

\newcommand{\Z}{\mathbb Z}
\newcommand{\Rm}{\mathcal R}
\newcommand{\I}{\mathcal I}
\newcommand{\SelCorr}{\mathrm{Sel}_{\mathrm{corr}}}
\newcommand{\SelStar}{\mathrm{Sel}^{\star}}
\newcommand{\eps}{\varepsilon}
\newcommand{\CJ}{\mathrm{carry\_jump}}
\newcommand{\WW}{\mathrm{wrap}}
\newcommand{\Flat}{\mathrm{flat}}
\newcommand{\OTHA}{\mathrm{other\_1000}}
\newcommand{\one}[1]{\mathbf 1_{\{#1\}}}

\title{Odd-modulus Hamilton decomposition of the directed $5$-torus\\
Further revised consolidated manuscript after the M1 front-end promotion review}
\author{}
\date{March 17, 2026}

\begin{document}
\maketitle

\begin{abstract}
This manuscript reorganizes the current odd-$m$, $d=5$ proof program after the
three theorem-side corrections made in notes 124--127.

First, the old theorem-side target ``future words of the current event map
separate states on the exact marked object'' is removed. For the actual current
event map $\eps_4$, that target is false. The correct theorem-side package is
instead the exact marked carry-jump-to-wrap object with anchor rigidity and the
comparison chain
\[
(B,\beta) \cong (\rho,\beta,q_0,\sigma) \twoheadrightarrow (\beta,\delta).
\]
Second, the intended theorem-side quotient is now fixed to
\[
Q_{\mathrm{th}}=(B,\beta),
\]
and this already closes the theorem-side M3 burden. Third, the historical
local/countdown-carrier provenance question is no longer a live proof target:
checked-range computation shows that the natural candidates
\[
Q_{\tau}=(B,\min(\tau,8),\eps_4),
\qquad
Q_{\mathrm{fw}}=(B,\eps_4,f_1,\dots,f_7)
\]
work at $m=11$ but fail from $m=13$ onward.

With those corrections in place, the proof program becomes cleaner. The
colorwise selector descent theorem already suffices for the generic upgrade, so
one no longer needs a single common selector family. The theorem-side part of
$M2$--$M3$ is now closed; $M4$ is closed on $\SelCorr$; the color-$4$ M5
branch is closed on $\SelStar$; and $M6$ remains closed. On the theorem-side
reduction front, the old monolithic $M1$ block now splits into two sharper
front-end promotions: the seed-isolation statement \textup{032}, and the
non-trigger part of the best-seed defect classification \textup{033}. The
exact trigger-family formula for $H_{L1}$ is already theorem-level through the
accepted \textup{033}/\textup{098} package. The remaining work is therefore
concentrated in those two front-end promotions and in the construction of a
full colorwise graph package for the remaining colors.

The purpose of this manuscript is to record that corrected proof spine in one
place and to state the remaining tasks as honest placeholders of the right
size.
\end{abstract}

\tableofcontents

\section{Introduction and module map}

Let
\[
D_5(m)=\operatorname{Cay}((\Z_m)^5,\{e_0,e_1,e_2,e_3,e_4\})
\]
be the directed $5$-torus. The long-term target is a decomposition of
$D_5(m)$ into five arc-disjoint directed Hamilton cycles for every odd modulus
$m$ in the D5 range.

The proof naturally splits into six packages, labelled $M1$ through $M6$.
After notes 124--127, their honest status is as follows.

\begin{center}
\renewcommand{\arraystretch}{1.15}
\begin{tabularx}{\textwidth}{@{}>{\bfseries}l>{\raggedright\arraybackslash}p{2.6cm}X@{}}
\toprule
Module & Status & Current content in this manuscript \\
\midrule
$M1$ & open, sharpened & reduced to front-end seed isolation \textup{032} and the non-trigger part of the best-seed defect classification \textup{033}; the exact $H_{L1}$ trigger-family formula is already imported \\
$M2$ & theorem-side closed in corrected form & exact marked carry-jump-to-wrap object with anchor rigidity and explicit coordinates equivalent to $(\rho,\beta,q_0,\sigma)$ \\
$M3$ & theorem-side closed & intended theorem-side quotient fixed to $(B,\beta)$; comparison/factor chain to $(\rho,\beta,q_0,\sigma)$ and then to $(\beta,\delta)$ \\
$M4$ & closed on $\SelCorr$ & symbolic all-odd corrected-selector theorem giving five global torus permutations \\
$M5$ & partially closed & the color-$4$ map of the surgery family $\SelStar$ is Hamilton for all odd $m$; the remaining colors are still open \\
$M6$ & closed & exceptional continuation / globalization package \\
\bottomrule
\end{tabularx}
\end{center}

The main conceptual corrections relative to the earlier consolidated draft are:
\begin{enumerate}[label=\textup{(\alph*)},leftmargin=2.5em]
\item the theorem-side M2/M3 package is no longer phrased using future-word
      separation for the current event map $\eps_4$;
\item the intended quotient is frozen theorem-side to $(B,\beta)$;
\item the historical local-provenance candidates $Q_{\tau}$ and
      $Q_{\mathrm{fw}}$ are treated as checked no-go objects, not as hidden
      targets of the final proof;
\item the graph-theoretic endgame uses a \emph{colorwise} selector package,
      so one does not need one common selector family before the final generic
      upgrade can be stated;
\item the old single $M1$ placeholder must be split: the exact trigger-family
      formula for $H_{L1}$ is already promoted by the accepted
      \textup{033}/\textup{098} package, while the genuine remaining front-end
      work is only seed isolation \textup{032} together with the non-trigger
      channel-classification part of \textup{033}.
\end{enumerate}

Accordingly, the current unconditional outputs are the following.

\begin{theorem}[Theorem-side M2/M3 package]
For odd $m$ in the accepted D5 regime, the exact marked carry-jump-to-wrap
object exists theorem-side with coordinates equivalent to
\[
(\rho,\beta,q_0,\sigma)
\qquad\text{and}\qquad
(B,\beta).
\]
The quotient $(B,\beta)$ is exact deterministic for the current event map
$\eps_4$, retains grouped base, and factors canonically to $(\beta,\delta)$.
\end{theorem}

\begin{theorem}[Corrected-selector M4 theorem for $\SelCorr$]
For every odd $m\ge 5$, the corrected selector family $\SelCorr$ defines five
global torus permutations whose arc sets are pairwise disjoint and cover the
full directed arc set of the torus.
\end{theorem}

\begin{theorem}[Color-$4$ M5 theorem for $\SelStar$]
For every odd $m\ge 5$, the color-$4$ torus map coming from the selector
surgery family $\SelStar$ is Hamilton.
\end{theorem}

\begin{theorem}[Globalization theorem]
For odd $m$ in the accepted D5 regime, the exceptional continuation /
globalization package is available and yields the standard raw exactness of
$(\beta,\delta)$ on the true accessible boundary union.
\end{theorem}

What remains open is therefore narrower than in the older manuscript: the two
front-end promotions \textup{032} and the non-trigger part of \textup{033},
and the missing graph-level colorwise exactness and Hamiltonicity packages for
the remaining colors.

\section{The corrected generic upgrade criterion}

The older consolidated draft still used two stronger hypotheses than the proof
actually needs:
\begin{enumerate}[label=\textup{(\roman*)},leftmargin=2.5em]
\item future-word separation on the exact marked object for the current event
      map $\eps_4$;
\item one common selector family carrying all five colors at once.
\end{enumerate}
Both are stronger than necessary. The corrected proof spine uses instead the
exact theorem-side anchor package and a colorwise selector package.

\subsection{Colorwise selector descent}

\begin{definition}[colorwise selector package]
Let
\[
V=(\Z/m\Z)^5
\]
be the vertex set of the directed $5$-torus.
A \emph{colorwise selector package} over $V$ consists of:
\begin{enumerate}[label=\textup{(C\arabic*)},leftmargin=2.9em]
\item a surjection $\pi:\Rm\to V$ from a representative space $\Rm$;
\item for each color $i\in\{0,1,2,3,4\}$, an invariant
      \[
      I_i:\Rm\to \I_i
      \]
      constant on every fiber of $\pi$;
\item for each color $i$, a local selected arc
      \[
      a_i(r)=(\pi(r),y_i(r))
      \]
      depending only on $I_i(r)$;
\item for each representative $r\in\Rm$, the five arcs
      \[
      a_0(r),a_1(r),a_2(r),a_3(r),a_4(r)
      \]
      are exactly the five outgoing torus arcs from $\pi(r)$;
\item for each color $i$ and each vertex $x\in V$, there is exactly one local
      color-$i$ arc ending at $x$.
\end{enumerate}
\end{definition}

\begin{theorem}[colorwise selector descent]
\label{thm:colorwise-selector-descent}
Assume a colorwise selector package over $V$.
For each color $i$, define
\[
F_i(x)=y_i(r)
\]
for any representative $r$ with $\pi(r)=x$.
Then:
\begin{enumerate}[label=\textup{(\alph*)},leftmargin=2.5em]
\item each $F_i$ is well defined;
\item each $F_i$ is a permutation of $V$;
\item the five arc sets
      \[
      A_i=\{(x,F_i(x)):x\in V\}
      \]
      are pairwise disjoint;
\item their union is the full directed arc set of the torus.
\end{enumerate}
\end{theorem}

\begin{proof}
Fix a color $i$ and a vertex $x\in V$.
If $r,r'\in\Rm$ satisfy $\pi(r)=\pi(r')=x$, then $I_i(r)=I_i(r')$ by fiber
constancy, so the color-$i$ selected arc agrees on $r$ and $r'$.
Thus $F_i(x)$ is well defined.

For a fixed representative $r$, axiom \textup{(C4)} says that the five color
arcs are exactly the five outgoing torus arcs from $\pi(r)$.
Hence the arc sets $A_0,\dots,A_4$ are pairwise disjoint and together cover the
full outgoing arc set.

Fix a color $i$.
Every vertex has out-degree $1$ in $A_i$ by definition.
By \textup{(C5)}, every vertex also has in-degree $1$ in $A_i$.
On a finite set this implies that $F_i$ is bijective, hence a permutation of
$V$.
\end{proof}

\subsection{Return-cycle lift}

\begin{theorem}[return-cycle lift theorem]
\label{thm:return-cycle-lift}
Let $F$ be a permutation of a finite set $V$.
Suppose there exist a subset $S\subseteq V$ and a positive integer $r$ such
that:
\begin{enumerate}[label=\textup{(R\arabic*)},leftmargin=2.8em]
\item every $F$-orbit meets $S$;
\item consecutive visits of an orbit to $S$ occur exactly $r$ steps apart;
\item the first-return map $T:S\to S$ induced by $F$ is a single cycle of
      length $|S|$;
\item every return segment of length $r$ contains exactly $r$ distinct
      vertices.
\end{enumerate}
Then $F$ is a single cycle of length $r|S|$.
\end{theorem}

\begin{proof}
Pick $s\in S$.
By (R2), the orbit of $s$ meets $S$ precisely at times $0,r,2r,\dots$.
By (R3), the successive return points
\[
s,T(s),T^2(s),\dots,T^{|S|-1}(s)
\]
are pairwise distinct and $T^{|S|}(s)=s$.
Equivalently,
\[
F^{kr}(s)\neq s\quad (0<k<|S|),
\qquad
F^{|S|r}(s)=s.
\]
By (R4), each return segment of length $r$ consists of $r$ distinct vertices,
and two distinct return segments cannot overlap because that would create an
earlier repeated point on the orbit. So the orbit of $s$ contains exactly
$r|S|$ distinct vertices.

Finally, by (R1) every orbit meets $S$, and by (R3) all points of $S$ already
lie on the orbit of $s$. Hence every orbit is the same orbit.
Therefore $F$ is a single cycle of length $r|S|$ on all of $V$.
\end{proof}

\begin{corollary}[iterated first-return reduction]
\label{cor:iterated-first-return}
Let $G$ be a permutation of a finite set $X$ and let $s:X\to\Z/m\Z$ satisfy
\[
s(Gx)=s(x)+1
\qquad\text{for all }x\in X.
\]
Set $Q:=\{x:s(x)=0\}$ and $H:=G^m|_Q$. Then $G$ is one cycle on $X$ if and
only if $H$ is one cycle on $Q$.
\end{corollary}

\begin{proof}
Apply Theorem~\ref{thm:return-cycle-lift} with $r=m$ and $S=Q$.
\end{proof}

\subsection{Corrected generic upgrade theorem}

\begin{theorem}[corrected generic D5 upgrade criterion]
\label{thm:generic-upgrade}
Let $m$ be odd.
Assume the following theorem packages.
\begin{enumerate}[label=\textup{(G\arabic*)},leftmargin=2.9em]
\item a channel-classification theorem reducing the full odd-$m$, $d=5$
      problem to the best-seed channel $R1\to H_{L1}$;
\item the corrected theorem-side M2/M3 package: an exact marked carry-jump-to-wrap
      object together with the intended exact deterministic quotient
      $Q_{\mathrm{th}}=(B,\beta)$;
\item a colorwise selector package on the full torus;
\item for each color $i\in\{0,1,2,3,4\}$, a return section $S_i$ satisfying
      the hypotheses of the return-cycle lift theorem for the global map $F_i$;
\item the exceptional continuation / globalization input already supplied by the
      accepted odd-$m$ D5 package.
\end{enumerate}
Then the directed torus $D_5(m)$ admits a decomposition into five arc-disjoint
Hamilton cycles.
\end{theorem}

\begin{proof}
By (G2), the canonical theorem-side coordinates already descend on the intended
quotient $Q_{\mathrm{th}}=(B,\beta)$.
By (G3) and Theorem~\ref{thm:colorwise-selector-descent}, the five local color
rules descend to well-defined global torus permutations
\[
F_0,F_1,F_2,F_3,F_4,
\]
whose arc sets are pairwise disjoint and cover the full directed arc set.
By (G4), each $F_i$ is Hamiltonian.
Therefore these five Hamilton cycles form a Hamilton decomposition of the
directed torus.
\end{proof}

\begin{remark}
This corrected generic theorem is weaker than the old one in exactly the right
ways: it no longer uses the false future-word-separation target for $\eps_4$,
and it no longer requires one common selector family before the final upgrade
can be stated.
\end{remark}

\section{The theorem-side M2/M3 package}

We summarize the corrected theorem-side package that replaces the obsolete
future-word formulation.

\begin{theorem}[corrected theorem-side M2 core]
\label{thm:corrected-M2-core}
Fix odd $m$ in the accepted D5 regime and the best-seed channel
\[
R1\to H_{L1}.
\]
There exists an exact marked carry-jump-to-wrap object $X_m$ with the following
properties:
\begin{enumerate}[label=\textup{(M2\arabic*)},leftmargin=2.9em]
\item $X_m$ admits equivalent exact coordinates
      \[
      (B,\beta)
      \qquad\text{and}\qquad
      (\rho,\beta,q_0,\sigma),
      \]
      where $B$ is grouped base, $\beta$ is the canonical clock, $\rho$ is the
      source residue, $q_0$ is the marked chain label, and $\sigma$ is the high
      digit;
\item the bare bridge quotient
      \[
      (\beta,q_0,\sigma)\cong(\beta,\delta),
      \qquad
      \delta=q_0+m\sigma,
      \]
      is exact deterministic for the current event map $\eps_4$;
\item on each regular fixed-$w$ slice, the return dynamics form the marked
      chain $q\mapsto q+1$ with unique carry state $q=m-1$;
\item the canonical clock $\beta$ descends to any exact deterministic quotient
      for $\eps_4$, and if such a quotient also retains grouped base $B$, then
      the full theorem-side readout $(B,\beta)$ descends.
\end{enumerate}
\end{theorem}

\begin{remark}
This is the corrected theorem-side replacement for the old sentence ``future
words of $\eps_4$ separate states on the exact marked object.'' The latter is
false; the anchor/marked-chain rigidity above is the theorem actually used
later.
\end{remark}

\begin{theorem}[theorem-side intended quotient verification]
\label{thm:Qth-verification}
The intended theorem-side quotient
\[
Q_{\mathrm{th}}=(B,\beta)
\]
is exact deterministic for the current event map $\eps_4$ and retains grouped
base. Consequently there is a canonical factor/comparison chain
\[
Q_{\mathrm{th}}
\twoheadrightarrow
(B,\beta)
\cong
(\rho,\beta,q_0,\sigma)
\twoheadrightarrow
(\beta,\delta).
\]
In particular, the theorem-side M3 burden is closed once the intended quotient
is fixed to $(B,\beta)$.
\end{theorem}

\begin{proof}
Grouped base retention is immediate because $B$ is a coordinate of
$Q_{\mathrm{th}}$.
Exactness of the current event map follows because $(B,\beta)$ determines the
enriched state $(\rho,\beta,q_0,\sigma)$, hence in particular the bare bridge
state $(\beta,q_0,\sigma)$, and the current event $\eps_4$ factors through that
bridge state.
Deterministic successor holds because $(B,\beta)$ is an exact coordinate system
on the exact marked object.
The displayed factor chain is therefore the theorem-side comparison theorem.
\end{proof}

\begin{proposition}[future-word separation for $\eps_4$ fails]
\label{prop:future-word-fails}
Future $\eps_4$-words do not separate states on the full exact marked object.
Already on the bridge quotient $(\beta,q_0,\sigma)$, the infinite future
$\eps_4$-word depends only on $(\beta,q_0)$ and is independent of $\sigma$.
Thus the bridge states split into exactly $m^2$ future-word classes of size
$m$.
\end{proposition}

\begin{remark}
This negative statement is not a side issue. It is the reason the old theorem-side
M2 wording had to be replaced.
\end{remark}

\section{The historical local-provenance no-go}

The theorem-side M3 package is now closed. A separate historical question is
whether older local/countdown-carrier packages can be identified canonically
with the theorem-side quotient $(B,\beta)$. The current checked answer is no.

\begin{proposition}[checked-range local-provenance no-go]
\label{prop:provenance-no-go}
Consider the two historical candidate quotients
\[
Q_{\tau}=(B,\min(\tau,8),\eps_4),
\qquad
Q_{\mathrm{fw}}=(B,\eps_4,f_1,\dots,f_7).
\]
On the regenerated exact marked carry-jump-to-wrap object, both candidates:
\begin{enumerate}[label=\textup{(\alph*)},leftmargin=2.5em]
\item satisfy the full intended-quotient test at $m=11$;
\item fail from $m=13$ onward;
\item from $m=13$ onward, fail both direct exact identification with
      $(B,\beta)$ and deterministic successor on the accessible quotient.
\end{enumerate}
In particular, the historical local-provenance claim is false for both tested
candidate families on the checked range
\[
m\in\{11,13,15,17,19,21,23,25\}.
\]
\end{proposition}

\begin{remark}[first exact obstruction at $m=13$]
The first obstruction already occurs at $m=13$.
There are two regular states with the same grouped base
\[
B=(4,1,4,10,\mathrm{regular})
\]
and the same local candidate state, but with different theorem-side clocks,
namely $\beta=9$ and $\beta=10$.
Their successors land in different next candidate states, so the historical
candidate quotient is not even deterministic on successor from that modulus
onward.
\end{remark}

\begin{remark}
Because of Proposition~\ref{prop:provenance-no-go}, the historical provenance
question is no longer part of the live proof-closing route. The final manuscript
should use $(B,\beta)$ theorem-side and should treat $Q_{\tau}$ and
$Q_{\mathrm{fw}}$ only as checked no-go comparison objects.
\end{remark}

\section{Closed graph-side branches}

\subsection{The all-odd M4 theorem for $\SelCorr$}

\begin{theorem}[M4 for $\SelCorr$]
\label{thm:M4-SelCorr}
For every odd $m\ge 5$, the corrected selector family $\SelCorr$ defines five
global torus permutations whose arc sets are pairwise disjoint and cover the
full directed arc set of the torus.
\end{theorem}

\begin{remark}
This is the fully symbolic all-odd selector theorem currently available on the
graph side. It closes one exact M4 package, but that family is not presently
an M5-friendly family.
\end{remark}

\subsection{The closed color-4 M5 branch for $\SelStar$}

\begin{theorem}[color-4 M5 for $\SelStar$]
\label{thm:M5-color4-SelStar}
For every odd $m\ge 5$, the color-$4$ map $F_4^{\star}$ of the selector surgery
family $\SelStar$ is Hamilton.
\end{theorem}

\begin{remark}
The proof uses the sharpened nested-return route:
\begin{enumerate}[label=\textup{(\arabic*)},leftmargin=2.5em]
\item an exact first-return clock on $P_0=\{\Sigma=0\}$;
\item an exact second nested clock on $P_1=\{\Sigma=0,\ x_3=0\}$;
\item a final $m^2$-state return map on
      \[
      P_2=\{\Sigma=0,\ x_3=0,\ x_0+x_2=0\};
      \]
\item identification of that true final return with the corrected-row model,
      whose row-stitching theorem gives one cycle.
\end{enumerate}
So the color-$4$ branch is a genuine theorem, not just checked evidence.
\end{remark}

\subsection{The imported globalization package}

\begin{proposition}[M6 remains closed]
\label{prop:M6-closed}
The odd-$m$ exceptional continuation / globalization package remains available
without change after the theorem-side corrections of notes 124--127.
\end{proposition}

\begin{proof}
The theorem-side corrections only change the bridge package used downstream.
They do not modify the accepted odd-$m$ exceptional continuation theorem.
\end{proof}

\section{Revised conditional final theorem and remaining placeholders}

With Sections 2--5 in place, the remaining proof burden is substantially smaller
than in the old consolidated draft.
One no longer needs:
\begin{enumerate}[label=\textup{(\roman*)},leftmargin=2.5em]
\item a theorem-side historical provenance identification for $Q_{\tau}$ or
      $Q_{\mathrm{fw}}$;
\item future-word separation for $\eps_4$;
\item a single common selector family before stating the final generic upgrade.
\end{enumerate}

The remaining tasks are now the following.
The M1 block has been sharpened by the present review: the exact $H_{L1}$
trigger-family formula is already promoted, so only the seed-isolation step and
the non-trigger channel-classification step remain external.

\begin{proposition}[Imported trigger-family theorem]
\label{prop:imported-HL1-trigger}
The exact hole-family formula
\[
H_{L1}=
\{(q,w,u,\Theta)=(m-1,m-1,u,2): u\neq 2\}
\]
already has a stable theorem-level home in the accepted \textup{033}/\textup{098}
package. It is not part of the remaining first-principles $M1$ burden of this
manuscript.
\end{proposition}

\begin{proof}
This is exactly the theorem content isolated in the promoted trigger-family note
\textup{033} and retained in the compact structural block \textup{098}.
\end{proof}

\begin{placeholder}[Front-end seed isolation / \textup{032}]
\label{ph:032}
Promote the best-endpoint-seed statement usually referred to as \textup{032}
into a manuscript-order theorem with explicit odd-$m$ quantifier: after the
earlier tiny static repair branches are pruned, the odd-$m$, $d=5$
endpoint-local front end reduces to the seed
\[
L=[2,2,1], \qquad R=[1,4,4].
\]
\end{placeholder}

\begin{placeholder}[Best-seed defect classification / non-trigger part of \textup{033}]
\label{ph:033-class}
For the seed $[2,2,1]/[1,4,4]$, promote the remaining channel-classification
content of the old \textup{033} analysis into one manuscript-order theorem: the
changed source families are exactly
\[
L3,\ R1,\ R2,\ R3,
\]
the channels $L3$, $R2$, and $R3$ repair directly, and the unique unresolved
channel is
\[
R1\to H_{L1}.
\]
The exact formula for $H_{L1}$ itself is already imported by
Proposition~\ref{prop:imported-HL1-trigger}.
\end{placeholder}

\begin{theorem}[M1 after the front-end promotions]
\label{thm:M1-from-032-and-033class}
Assume Placeholders~\ref{ph:032} and~\ref{ph:033-class}.
Then every odd-$m$, $d=5$ instance is either already closed by an earlier
front-end repair branch or reduces to the best-seed channel
\[
R1\to H_{L1}.
\]
Moreover the target hole family is exactly
\[
H_{L1}=
\{(q,w,u,\Theta)=(m-1,m-1,u,2): u\neq 2\}.
\]
\end{theorem}

\begin{proof}
Placeholder~\ref{ph:032} reduces the endpoint-local front end to the best seed
$[2,2,1]/[1,4,4]$.
Placeholder~\ref{ph:033-class} says that inside that seed the only unresolved
channel is $R1\to H_{L1}$ after the direct repairs $L3$, $R2$, and $R3$ are
removed.
Proposition~\ref{prop:imported-HL1-trigger} supplies the exact formula for the
target hole family.
Therefore the full odd-$m$, $d=5$ problem reduces to the stated best-seed
channel.
\end{proof}

\begin{remark}[Promotion review of \textup{032} and \textup{033}]
The present review does \emph{not} justify promoting the whole
\textup{032}/\textup{033} front-end package without further work.
In the supplied bundles, the exact trigger-family formula for $H_{L1}$ is
already theorem-level, but the seed-isolation statement \textup{032} and the
broader direct-repair / unique-unresolved-channel part of \textup{033} are
still front-end artifact or frontier statements rather than integrated
manuscript-order proofs.
So the correct replacement for the old single $M1$ placeholder is the pair of
Placeholders~\ref{ph:032} and~\ref{ph:033-class}, not a claim that all of $M1$
is now unconditional.
\end{remark}

\begin{placeholder}[Colorwise graph package]
\label{ph:colorwise-graph}
Produce five local color rules
\[
a_0,a_1,a_2,a_3,a_4
\]
on the full torus which together form a colorwise selector package in the sense
of Section~2.
The proof no longer needs these five rules to arise from one common
permutation-valued selector family.
It only needs fiberwise well-defined local arcs, pointwise outgoing
exhaustiveness, and incoming uniqueness for each color.
\end{placeholder}

\begin{placeholder}[Remaining colorwise M5 package]
\label{ph:remaining-M5}
For the global maps supplied by Placeholder~\ref{ph:colorwise-graph}, prove
Hamiltonicity for the colors not already covered by the closed
$F_4^{\star}$ theorem. In practice this means either:
\begin{enumerate}[label=\textup{(\alph*)},leftmargin=2.5em]
\item extend the nested-return / monodromy analysis to the additional selected
      colors; or
\item replace the current color rules by a compatible colorwise family whose
      remaining maps admit a cleaner return-cycle theorem.
\end{enumerate}
\end{placeholder}

\begin{theorem}[revised conditional final theorem]
\label{thm:revised-conditional-final}
Assume odd $m$ lies in the accepted D5 regime and that
Placeholders~\ref{ph:032}, \ref{ph:033-class}, \ref{ph:colorwise-graph},
and \ref{ph:remaining-M5} have been resolved.
Then the directed torus
\[
D_5(m)=\operatorname{Cay}((\Z_m)^5,\{e_0,e_1,e_2,e_3,e_4\})
\]
admits a decomposition into five arc-disjoint directed Hamilton cycles.
\end{theorem}

\begin{proof}
Theorem~\ref{thm:M1-from-032-and-033class} supplies the channel reduction
package once the front-end promotions
Placeholders~\ref{ph:032} and~\ref{ph:033-class} are resolved.
The corrected theorem-side M2/M3 package is already available by
Theorems~\ref{thm:corrected-M2-core} and \ref{thm:Qth-verification}.
Placeholder~\ref{ph:colorwise-graph} supplies the colorwise selector package.
Placeholder~\ref{ph:remaining-M5} supplies the missing return-cycle theorems.
Proposition~\ref{prop:M6-closed} supplies the globalization input.
Therefore Theorem~\ref{thm:generic-upgrade} applies and yields a Hamilton
decomposition.
\end{proof}

\section{Next order of work}

After the local-provenance no-go, the next order is cleaner than before.

\begin{enumerate}[label=\textup{(\arabic*)},leftmargin=2.7em]
\item \textbf{Freeze the theorem-side package.}
      Treat Sections 3 and 4 as stable.
      In particular, do not spend further proof effort trying to identify
      $Q_{\tau}$ or $Q_{\mathrm{fw}}$ canonically with $(B,\beta)$.

\item \textbf{Promote \textup{032} and the non-trigger part of \textup{033}.}
      The exact $H_{L1}$ trigger formula is already theorem-level; the genuine
      remaining theorem-side reduction burden is to lift the seed-isolation and
      the direct-repair / unique-unresolved-channel facts into manuscript order.

\item \textbf{Use the colorwise graph formulation, not family unification.}
      The graph endgame should be organized around Placeholder
      \ref{ph:colorwise-graph}, not around the stronger and unnecessary demand
      for one common selector family.

\item \textbf{Pursue the next M5 branch with the strongest current evidence.}
      After the closed color-$4$ branch, the next natural candidate is the
      color-$3$ surgery branch, since the existing probes already suggest that
      it is structurally closer to the color-$4$ route than colors $0,1,2$.

\item \textbf{Only after one more colorwise Hamiltonicity theorem is in hand,
      decide whether to continue branch-by-branch or redesign the colorwise
      package.}
      At that stage the graph package will be concrete enough to show whether a
      mixed colorwise assembly is realistic or whether a new surgery family is
      needed.
\end{enumerate}

\begin{remark}[bottom line]
The proof-closing picture after 127 is simpler than before.
The theorem-side M2/M3 package is now closed in corrected form, the historical
local provenance claim is checked false on the currently tested range, and the
remaining open work sits in the front-end promotions \textup{032} and the
non-trigger part of \textup{033}, and in the graph-side colorwise package for
the remaining colors.
That is the right corrected integrated manuscript state for the next closure
steps.
\end{remark}

\end{document}
